% =========================================================
\section{Adiabatic BLIMP-Inspired Skin Friction Model}
\label{sec:blimp_cebeci_smith}
% =========================================================

\subsection{Scope and modelling hierarchy}

The nozzle boundary layer is evaluated with a reduced-order procedure inspired by
the JANNAF BLIMP-J methodology (NASA-CR-144046) \cite{BLIMPJ}. The implementation is referred to
here as \emph{BLIMP-lite}: it retains a wall-normal profile, axisymmetric continuity,
the compressible thin-layer momentum equation and the Cebeci--Smith turbulent
closure, but it does not reproduce BLIMP-J's chemically reacting energy and species
equations and is not a reimplementation of BLIMP-J. It is intended for attached-flow preliminary design and geometry
optimization at substantially lower cost than a RANS calculation.

At each wall station, the quasi-one-dimensional core-flow model supplies the edge
quantities

\begin{equation}
U_e(s),\quad p_e(s),\quad T_e(s),\quad \rho_e(s),\quad
M_e(s),\quad \gamma_e(s),\quad \mu_e(s),\quad Pr_e(s),
\end{equation}

where \(s\) is distance along the wall and \(y\) is the inward wall-normal
coordinate. RocketCEA chamber, throat and exit properties are interpolated locally.
The wall contour supplies the axisymmetric radius \(r(s,y)\).

\subsection{Adiabatic wall condition}

Wall temperature is not prescribed as an arbitrary fixed input. Without a conjugate
wall-conduction and coolant model, a fixed gas-side value cannot be predicted
independently of the final geometry. The present gas-side calculation therefore
adopts an adiabatic wall,

\begin{equation}
\boxed{T_w(s)=T_r(s)=T_{aw}(s)},
\qquad
\boxed{q_w(s)=0}.
\label{eq:blimp_adiabatic_wall}
\end{equation}

For the fully turbulent closure, the local recovery factor and adiabatic-wall
temperature are

\begin{equation}
r_f(s)=Pr_e(s)^{1/3},
\label{eq:blimp_recovery_factor}
\end{equation}

\begin{equation}
\boxed{
T_{aw}(s)=T_e(s)
\left[
1+r_f(s)\frac{\gamma_e(s)-1}{2}M_e(s)^2
\right]
}.
\label{eq:blimp_recovery_temperature}
\end{equation}

The wall-normal temperature profile is closed using the adiabatic Walz relation,

\begin{equation}
\boxed{
T(s,y)=T_{aw}(s)-\left[T_{aw}(s)-T_e(s)\right]
\left(\frac{u(s,y)}{U_e(s)}\right)^2
}.
\label{eq:blimp_walz_adiabatic}
\end{equation}

Equation~\eqref{eq:blimp_walz_adiabatic} satisfies \(T=T_{aw}\) and
\(\partial T/\partial y=0\) at the wall, and \(T=T_e\) at the boundary-layer
edge. Density is evaluated locally from

\begin{equation}
\rho(s,y)=\frac{p_e(s)}{R(s)T(s,y)},
\label{eq:blimp_density}
\end{equation}

while molecular viscosity is obtained by log-interpolation of the RocketCEA
transport values, with a mild temperature power-law extrapolation outside their
tabulated interval.

\subsection{Axisymmetric compressible boundary-layer equations}

The thin-layer continuity equation is

\begin{equation}
\boxed{
\frac{\partial}{\partial s}\!\left(\rho u r\right)
+
\frac{\partial}{\partial y}\!\left(\rho v r\right)=0
},
\label{eq:blimp_continuity}
\end{equation}

and is integrated from the impermeable wall to obtain the normal velocity \(v\),
with \(v_w=0\). The streamwise momentum equation is

\begin{equation}
\boxed{
\rho u\frac{\partial u}{\partial s}
+\rho v\frac{\partial u}{\partial y}
=-\frac{dp_e}{ds}
+\frac{1}{r}\frac{\partial}{\partial y}
\left[
r\left(\mu+\mu_t\right)\frac{\partial u}{\partial y}
\right]
}.
\label{eq:blimp_momentum}
\end{equation}

The inviscid edge solution supplies the pressure-gradient forcing, equivalently

\begin{equation}
-\frac{dp_e}{ds}=\rho_eU_e\frac{dU_e}{ds}.
\label{eq:blimp_edge_pressure_gradient}
\end{equation}

The boundary conditions are

\begin{equation}
u(s,0)=0,\qquad v(s,0)=0,\qquad
u(s,y_e)=U_e(s),\qquad T(s,y_e)=T_e(s).
\label{eq:blimp_boundary_conditions}
\end{equation}

\subsection{Cebeci--Smith turbulent closure}

Reynolds shear is represented through the Boussinesq eddy-viscosity hypothesis,

\begin{equation}
-\rho\overline{u'v'}=\mu_t\frac{\partial u}{\partial y}.
\label{eq:blimp_boussinesq}
\end{equation}

The inner Cebeci--Smith layer uses a Van-Driest-damped mixing length
\cite{CebeciSmith,BLIMPVerification},

\begin{equation}
\ell=\kappa y
\left[1-\exp\left(-\frac{y^+}{A^+}\right)\right],
\qquad
\kappa=0.40,\quad A^+=26,
\label{eq:blimp_mixing_length}
\end{equation}

where

\begin{equation}
y^+=\frac{\rho_wu_\tau y}{\mu_w},
\qquad
u_\tau=\sqrt{\frac{\tau_w}{\rho_w}}.
\label{eq:blimp_wall_units}
\end{equation}

The corresponding inner eddy viscosity is

\begin{equation}
\mu_{t,i}=\rho\ell^2
\left|\frac{\partial u}{\partial y}\right|.
\label{eq:blimp_inner_eddy_viscosity}
\end{equation}

The outer layer is represented by the displacement-thickness wake scale

\begin{equation}
\mu_{t,o}
=0.0168\,\rho_eU_e\delta^*F_K(y),
\label{eq:blimp_outer_eddy_viscosity}
\end{equation}

with the Klebanoff intermittency function

\begin{equation}
F_K(y)=
\left[
1+5.5\left(\frac{y}{6\delta^*}\right)^6
\right]^{-1}.
\label{eq:blimp_klebanoff}
\end{equation}

The two regions are joined algebraically according to

\begin{equation}
\boxed{\mu_t=\min\!\left(\mu_{t,i},\mu_{t,o}\right)}.
\label{eq:blimp_eddy_join}
\end{equation}

Van Driest damping ensures \(\mu_t\rightarrow0\) at the wall. Consequently, wall
shear is obtained from the resolved molecular gradient rather than from an imposed
flat-plate skin-friction correlation:

\begin{equation}
\boxed{
\tau_w=\mu_w\left.\frac{\partial u}{\partial y}\right|_w
},
\qquad
\boxed{
C_f=\frac{2\tau_w}{\rho_eU_e^2}
}.
\label{eq:blimp_skin_friction}
\end{equation}

Since \(y^+\), \(\mu_t\), the temperature-dependent properties and \(\tau_w\) are
mutually dependent, they are updated iteratively at every streamwise station.

\subsection{Integral thicknesses and viscous losses}

The density-weighted compressible displacement and momentum thicknesses are

\begin{equation}
\boxed{
\delta^*(s)=\int_0^{y_e}
\left(1-\frac{\rho u}{\rho_eU_e}\right)dy
},
\label{eq:blimp_displacement_thickness}
\end{equation}

\begin{equation}
\boxed{
\theta(s)=\int_0^{y_e}
\frac{\rho u}{\rho_eU_e}
\left(1-\frac{u}{U_e}\right)dy
},
\label{eq:blimp_momentum_thickness}
\end{equation}

and \(H=\delta^*/\theta\). Displacement blockage is introduced through

\begin{equation}
r_{eff}(s)=r_w(s)-\delta^*(s),
\qquad
A_{eff}(s)=\pi r_{eff}(s)^2.
\label{eq:blimp_effective_area}
\end{equation}

Skin friction is retained as a distinct axial force. Since
\(dA_w=2\pi r\,ds\) and the axial tangent projection is \(dx/ds\),

\begin{equation}
\boxed{
F_{fric}=\int_{x_c}^{x_e}2\pi r_w(x)\tau_w(x)\,dx
},
\qquad
\boxed{
C_{F,fric}=\frac{F_{fric}}{p_cA_t}
}.
\label{eq:blimp_friction_force}
\end{equation}



Blockage and wall shear are thus represented separately and are not hidden inside a
single empirical efficiency.

\subsection{Numerical procedure and limitations}

The equations are marched implicitly over 61 streamwise stations. Each station uses
41 wall-normal nodes clustered toward the wall. The continuity equation provides
\(v(y)\); the momentum equation is discretized in conservative radial form; and
Picard iterations update \(u\), \(T\), \(\rho\), \(\mu\), \(y^+\), \(\mu_t\) and
\(\tau_w\). The converged integral quantities are interpolated onto the complete
geometry grid. In a default-case grid comparison against 91 streamwise by 57
wall-normal nodes, exit \(C_f\) and integrated \(C_{F,fric}\) changed by less than
0.5\%, while exit \(\delta^*\) changed by approximately 4\%.

The model assumes an attached, fully turbulent, axisymmetric boundary layer with no
wall mass transfer. The inlet profile is an engineering initialization because the
upstream chamber boundary-layer history is not otherwise available. Transition,
relaminarization, shock--boundary-layer interaction and separated flow are not
predicted. The adiabatic Walz relation \cite{Walz} replaces the full BLIMP-J reacting energy and
species system. Cebeci--Smith is an algebraic equilibrium closure; final designs
therefore still require mesh-refined RANS/CFD and experimental validation.
